% Characterization chapter.

\section{Factorisation through ego-net types}\label{sec:characterization}

Theorem~\ref{thm:dichotomy} places every MPNN output in the class of continuous
$1$-WL invariant functionals. This chapter gives an explicit form: the
output of a depth-$T$ MPNN on a featured graph is a continuous function
of the empirical distribution of the graph's rooted $T$-hop ego-net
types. The factorisation gives an approximation procedure for any
trained MPNN estimator that uses only local per-vertex statistics of
the graph.

\subsection{Rooted ego-nets}

Let $(\gph, x)$ be a featured graph and let $v \in \vertices$. For an
integer $T \ge 1$, define the \emph{$T$-hop ball around $v$} as the
set $B_T(v) = \{u \in \vertices : d_\gph(u, v) \le T\}$, where
$d_\gph$ denotes graph distance in $\gph$.

\begin{definition}[Rooted ego-net]\label{def:egonet}
The \emph{rooted $T$-hop ego-net at $v$}, written $\mathcal{E}_T(v)$,
is the triple
\[
\mathcal{E}_T(v) = (\gph[B_T(v)],\ x|_{B_T(v)},\ v),
\]
where $\gph[B_T(v)]$ is the subgraph of $\gph$ induced on $B_T(v)$,
$x|_{B_T(v)}$ is the restriction of the feature $x$ to $B_T(v)$, and
the final component marks $v$ as the root.
\end{definition}

Two rooted ego-nets $(\gph_1, x_1, v_1)$ and $(\gph_2, x_2, v_2)$ are
\emph{rooted-isomorphic} if there is an isomorphism of featured graphs
$\pi : (\gph_1, x_1) \to (\gph_2, x_2)$ with $\pi(v_1) = v_2$: same
edge structure, same features, root maps to root. Write $\mathcal{R}_T$
for the set of rooted-isomorphism classes of possible $T$-hop
ego-nets. Since the ego-net has at most $\sum_{k=0}^T d^k$ vertices in
a graph of maximum degree $d$, $\mathcal{R}_T$ is finite for every
fixed $T$ and $d$.

\begin{definition}[Ego-net type distribution]\label{def:pi-t}
Let $\pi_T(\gph, x) \in \Delta(\mathcal{R}_T)$ denote the empirical
distribution of ego-net types across the vertices of $\gph$:
\[
\pi_T(\gph, x)(\rho) = \frac{|\{v \in \vertices : [\mathcal{E}_T(v)] =
\rho\}|}{|\vertices|}
\qquad \text{for } \rho \in \mathcal{R}_T,
\]
where $[\mathcal{E}_T(v)]$ denotes the rooted-isomorphism class of
$\mathcal{E}_T(v)$.
\end{definition}

Rooted ego-nets and $1$-WL colour refinement are related: the
depth-$T$ colour $c_v^{(T)}$ of $v$ under Definition~\ref{def:wl} is a canonical
name for the rooted-isomorphism class of $\mathcal{E}_T(v)$. Colour
refinement is precisely the standard algorithm for computing this
canonical name.

\begin{example}\label{ex:egonet-c6}
Take the six-cycle $B = C_6$ with the alternating feature of
Example~\ref{ex:wl-feat} (odd vertices $+1$, even vertices $-1$). Compute
$\mathcal{E}_1(1)$: the $1$-hop ball is $B_1(1) = \{1, 2, 6\}$; the
induced subgraph has two edges ($\{1, 2\}$ and $\{1, 6\}$); the
features are $x(1) = +1$, $x(2) = -1$, $x(6) = -1$; the root is $1$.
Every odd vertex of $B$ has an $\mathcal{E}_1$-type
rooted-isomorphic to this one. Every even vertex has an
$\mathcal{E}_1$-type with roles reversed. The distribution $\pi_1(B,
x)$ places mass $1/2$ on each of these two types.
\end{example}

\subsection{Factorisation}

\begin{proposition}\label{prop:egonet-factorisation}
Let $M \in \mpnn$ be a message-passing GNN of depth $T$ with
continuous update and readout. There exists a continuous function
$F_M : \Delta(\mathcal{R}_T) \to \R$ such that
\[
M(\gph, x) = F_M\!\left(\pi_T(\gph, x)\right)
\]
for every featured graph $(\gph, x)$.
\end{proposition}

\begin{proof}
Fix $M \in \mpnn$ of depth $T$. By the proof of Proposition~\ref{prop:upper},
$M(\gph, x)$ is a permutation-invariant function of the multiset
$\{\!\{h_v^{(T)} : v \in \vertices\}\!\}$. The representation
$h_v^{(T)}$ is a function of the featured rooted $T$-hop ego-net
$\mathcal{E}_T(v)$ alone (since a depth-$T$ MPNN aggregates information
from the $T$-hop neighbourhood of each vertex). Two vertices with the
same rooted-isomorphism class of $\mathcal{E}_T$ receive the same
representation $h_v^{(T)}$.

Consequently, $\{\!\{h_v^{(T)} : v \in \vertices\}\!\}$ is determined
by the empirical distribution $\pi_T(\gph, x)$: for each type $\rho
\in \mathcal{R}_T$ occurring at some vertex, one instance of the
corresponding representation appears with multiplicity $|\vertices|
\cdot \pi_T(\gph, x)(\rho)$.

The readout, being permutation-invariant and continuous, defines a
continuous $F_M$ on $\Delta(\mathcal{R}_T)$ with $M(\gph, x) =
F_M(\pi_T(\gph, x))$.
\end{proof}

\subsection{Approximation from ego-net statistics}

Proposition~\ref{prop:egonet-factorisation} lets us approximate any MPNN
estimator's output by polynomial functions of the coordinates of
$\pi_T$ on a compact featured-graph class.

\begin{proposition}\label{prop:reproduce}
Let $M \in \mpnn$ have depth $T$, continuous update, and continuous
readout. Fix a compact featured-graph class $\mathcal{G}$ of graphs
on at most $n$ vertices with maximum degree at most $d$. For every
$\varepsilon > 0$ there exists a polynomial $P_\varepsilon$ in the
coordinates of $\pi_T$ such that
\[
\sup_{(\gph, x) \in \mathcal{G}} \left|M(\gph, x) - P_\varepsilon(\pi_T(\gph, x))\right|
\le \varepsilon.
\]
\end{proposition}

\begin{proof}
By Proposition~\ref{prop:egonet-factorisation}, $M(\gph, x) = F_M(\pi_T(\gph,
x))$ for a continuous $F_M : \Delta(\mathcal{R}_T) \to \R$. The
image of $\mathcal{G}$ under $\pi_T$ is a compact subset of
$\Delta(\mathcal{R}_T)$, itself a compact subset of $\R^{|\mathcal{R}_T|}$;
$\mathcal{R}_T$ is finite because a rooted $T$-hop ego-net has at
most $\sum_{k=0}^T d^k$ vertices. By the Stone-Weierstrass theorem,
continuous functions on this compact set are uniformly approximable
by polynomials in the coordinates.
\end{proof}

\subsection{The graph-blind statistic}\label{sec:graph-blind}

At $T = 1$ the ego-net type of $v$ is determined by $x_v$, $\deg(v)$,
and the multiset $\{\!\{x_u : u \in \neigh{v}\}\!\}$. Any statistic
that depends on the featured graph only through the empirical
distribution of $T = 1$ ego-net types is what we call a
\emph{graph-blind statistic}: it treats the graph as a bag of
per-vertex tuples $(x_v, \deg(v), \{\!\{x_u : u \in \neigh{v}\}\!\})$
and ignores everything else.

\begin{definition}[Graph-blind statistic]\label{def:graph-blind}
A functional $S : \mathcal{G} \to \R$ is \emph{graph-blind at depth
$1$} if $S(\gph, x) = S(\gph', x')$ whenever $\pi_1(\gph, x) =
\pi_1(\gph', x')$.
\end{definition}

Any graph-blind statistic is a function of the empirical distribution
of first-order neighbourhood tuples. Standard examples include edge
homophily and Newman attribute assortativity. Any ridge regression
built from vertex features constructed as functions of $(x_v, \deg(v),
\{\!\{x_u : u \in \neigh{v}\}\!\})$ is a graph-blind statistic.

Proposition~\ref{prop:reproduce} guarantees that a polynomial in $\pi_T$
approximates any continuous MPNN output to arbitrary precision. A
polynomial in $\pi_1$ can only approximate an MPNN whose factorisation
$F_M$ depends on $\pi_T$ only through its marginalisation to $\pi_1$;
this is a stronger property, and not every MPNN satisfies it.

The empirical audit of Section~\ref{sec:emp-reddit} tests this stronger
property directly on EchoGAE: it fits a ridge regression on eight
per-vertex features (all functions of the depth-$1$ ego-net type) to
EchoGAE's per-node output and reports the cross-validated $R^2$. On
the abortion dataset the $R^2$ is $0.60$; on the gun dataset the
$R^2$ is $-0.01$. A high $R^2$ is evidence that EchoGAE's
factorisation depends on $\pi_T$ only through its depth-$1$
marginalisation. A low $R^2$ is not evidence to the contrary: it
means only that the specific eight features chosen do not span the
required subspace of $\pi_1$.

\subsection{The audit procedure}

The audit procedure derived from this chapter has three steps.

\emph{Fit a graph-blind regression.} Choose a small vector of
per-vertex features, all functions of $(x_v, \deg(v), \{\!\{x_u : u
\in \neigh{v}\}\!\})$. Fit a ridge regression from these features to
the estimator's per-node output on the training distribution in
cross-validation. Report the cross-validated $R^2$.

\emph{Apply a 1-WL preserving rewiring.} Apply a graph rewiring that
preserves every $1$-WL colour class at every depth. The
degree-and-block-preserving edge-swap procedure of
Section~\ref{sec:emp-reddit} approximates this: it preserves the sign-block
edge multiplicities exactly and preserves depth-$1$ WL colours
approximately. Recompute the estimator's output on the rewired graph.

\emph{Compare cross-topic scores.} Compare the estimator's numerical
output on multiple datasets. If the estimator returns
indistinguishable scores on datasets the authors label as polarised
and non-polarised, the estimator is not discriminating between the
two categories.

Each check falls out of one of the results in this chapter or of
Theorem~\ref{thm:dichotomy}. The audit in the next chapter runs all three
on EchoGAE.
